\begin{tikzpicture}[
    node distance=0.4cm,
    box/.style={draw, rounded corners=2pt, fill=#1, text=white, font=\footnotesize\bfseries, align=left, inner sep=4pt, minimum width=8cm},
    box/.default=blue!60!black,
    optbox/.style={box=orange!80!black},
    evalbox/.style={box=green!50!black},
    modelbox/.style={box=purple!60!black},
    smallbox/.style={draw, rounded corners=2pt, fill=blue!15, font=\scriptsize, align=left, inner sep=3pt},
    arr/.style={->, >=stealth, thick, color=black!70},
    label/.style={font=\scriptsize\itshape, color=black!60},
]

% DATA SOURCES
\node[box] (data) {
\textbf{DATA SOURCES}\\[2pt]
SemEval ASTE: Rest14 (train), Rest14/15/16/Laptop14 (test)\\
DMASTE: 4 source domains (electronics, beauty, fashion, home) → 4 target (book, grocery, pet, toy)
};

% CANONICAL
\node[box=black!70, below=of data, minimum width=8cm] (canonical) {
\textbf{CANONICAL FORMAT}\\[1pt]
\{sentence, tokens, annotations[\{aspect, opinion, polarity\}]\}
};
\draw[arr] (data) -- (canonical);

% AUGMENTATION
\node[optbox, below=of canonical, minimum width=8cm] (aug) {
\textbf{AUGMENTATION} (optional)\\[1pt]
├ Aspect masking (configurable fraction, replace with <extra\_id>)\\
├ Duplication (controlled data multiplication)\\
└ Cross-domain data mixing (DMASTE fractions added to Rest14)
};
\draw[arr] (canonical) -- (aug);

% SYNTAX
\node[optbox, below=of aug, minimum width=8cm] (syntax) {
\textbf{SYNTAX ENRICHMENT} (optional)\\[1pt]
├ Compact dependency (separate Syntax: line, content words only)\\
├ Compact POS (POS tags for content words)\\
└ Inline (dep. or POS annotations mixed into input tokens)
};
\draw[arr] (aug) -- (syntax);

% FORMAT CONVERSION
\node[box, below=of syntax, minimum width=8cm] (format) {
\textbf{GENERATIVE FORMAT CONVERSION}\\[1pt]
├ Task subset: full triplet | singles | pairs\\
│\quad └ Curriculum: waypoint interpolation across epochs\\
├ Output format:\\
│\quad ├ Structured: [aspect, opinion, polarity]\\
│\quad ├ NL: ``\{aspect\} is described as \{opinion\}, expressing a \{polarity\} sentiment''\\
│\quad └ MvP markers: [A] x [O] y [S] z [SSEP]\\
├ Multi-ordering: 6 permutations of element order\\
└ STAR: pairwise sub-task decomposition (7× data)
};
\draw[arr] (syntax) -- (format);

% PER-EPOCH
\node[smallbox, below=0.3cm of format] (epoch) {Per-epoch re-randomisation (all above re-executed each epoch with different seed)};
\draw[arr] (format) -- (epoch);

% MODELS
\node[box, below=0.3cm of epoch, minimum width=5.5cm, xshift=-1.5cm] (t5) {
\textbf{FLAN-T5-base} (250M)\\[1pt]
AdamW, lr=3e-4, cosine, 30 epochs\\
Label smoothing=0.05, greedy decoding
};
\node[modelbox, below=0.3cm of epoch, minimum width=3.8cm, xshift=3.2cm] (roberta) {
\textbf{RoBERTa BASELINE} (125M)\\[1pt]
BIO tagger + biaffine scorer\\
+ polarity classifier
};
\draw[arr] (epoch.south) -- ++(0,-0.15) -| (t5.north);
\draw[arr] (epoch.south) -- ++(0,-0.15) -| (roberta.north);

% PARSING
\node[box=black!70, below=0.6cm of t5.south east, anchor=north, minimum width=8cm, xshift=0.5cm] (parse) {
\textbf{OUTPUT PARSING}\\[1pt]
├ NL: regex template matching\\
├ Structured: bracket/comma splitting\\
├ MvP: [SSEP] separator + marker extraction\\
└ Malformed outputs → score 0 (strict, no partial credit)
};
\draw[arr] (t5.south) -- ++(0,-0.3) -| (parse.north west);
\draw[arr] (roberta.south) -- ++(0,-0.3) -| (parse.north east);

% EVALUATION
\node[evalbox, below=of parse, minimum width=8cm] (eval) {
\textbf{EVALUATION}\\[1pt]
├ Exact-match micro P/R/F1 (primary metric)\\
├ Macro P/R/F1\\
├ Lenient F1 (token overlap, threshold=0.8)\\
├ Lenient F1 (longest common subsequence, threshold=0.8)\\
├ Soft F1 (embedding cosine similarity, aspect key, all-MiniLM-L6-v2)\\
├ Average token overlap + LCS overlap + match rate\\
└ Multi-seed protocol: 5 seeds (42, 123, 456, 789, 1337), mean ± std
};
\draw[arr] (parse) -- (eval);

\end{tikzpicture}
